\booksection{Apêndice do Capítulo 11: Conhecimento Técnico de Imagem Profissional}

\chapterintro{Para entender verdadeiramente os segredos da fotografia e do vídeo com celular, você precisa conhecer alguns fundamentos técnicos de imagem profissional. Esse conhecimento não só ajuda a usar melhor a câmera do celular, como também lhe dá mais critério na hora de escolher um aparelho.}

\booksubsection{I. Resolução, Taxa de Quadros e Bitrate: Os Três Pilares da Qualidade de Vídeo}

\subsubsection*{Resolução (Resolution):}

A resolução determina a nitidez da imagem, expressa em número de pixels.

\begin{itemize}
  \item \textbf{720p (HD)}: 1280×720 pixels, nitidez de nível básico.
  \item \textbf{1080p (Full HD)}: 1920×1080 pixels, a resolução mais popular atualmente.
  \item \textbf{2K (QHD)}: 2560×1440 pixels, mais nítido que 1080p.
  \item \textbf{4K (UHD)}: 3840×2160 pixels, ultra-alta definição, rico em detalhes.
  \item \textbf{8K}: 7680×4320 pixels, padrão do futuro, suporte atual limitado em dispositivos.
\end{itemize}

Quanto maior a resolução, mais nítida a imagem, mas também maior o arquivo e maior a exigência de armazenamento e processamento.

\subsubsection*{Taxa de Quadros (Frame Rate):}

A taxa de quadros determina a fluidez da imagem, expressa em quadros por segundo (fps).

\begin{itemize}
  \item \textbf{24 fps}: taxa de quadros padrão do cinema, com uma "sensação cinematográfica" única.
  \item \textbf{30 fps}: comum em TV e vídeos online, fluido e natural.
  \item \textbf{60 fps}: alta taxa de quadros, ideal para cenas de esporte e gravação de jogos, muito fluido.
  \item \textbf{120 fps / 240 fps}: taxa de quadros ultra-alta, usada para vídeos em câmera lenta.
\end{itemize}

Quanto maior a taxa de quadros, mais fluida a imagem e melhor o efeito de câmera lenta, mas também maior o arquivo.

\subsubsection*{Bitrate:}

O bitrate determina a densidade de dados do vídeo, expresso em quantidade de dados por segundo (Mbps).

\begin{itemize}
  \item Quanto maior o bitrate, melhor a qualidade de imagem, mais detalhes preservados, mas também maior o arquivo.
  \item Bitrate muito baixo causa artefatos de compressão, como blocos e manchas na imagem.
  \item A gravação de vídeo profissional geralmente exige alto bitrate para preservar mais margem para pós-produção.
\end{itemize}

\booksubsection{II. Formatos de Codificação de Vídeo: H.264, H.265 e ProRes}

\subsubsection*{H.264 (AVC):}

\begin{itemize}
  \item \textbf{Característica}: formato de codificação mais amplamente utilizado, compatibilidade excelente.
  \item \textbf{Vantagens}: quase todos os dispositivos conseguem reproduzir, eficiência de compressão moderada.
  \item \textbf{Desvantagens}: comparado a novos formatos, eficiência de compressão menor, arquivos maiores.
  \item \textbf{Cenário de uso}: compartilhamento de vídeos online, gravação cotidiana do usuário comum.
\end{itemize}

\subsubsection*{H.265 (HEVC):}

\begin{itemize}
  \item \textbf{Característica}: versão aprimorada do H.264, eficiência de compressão cerca de 50\% maior.
  \item \textbf{Vantagens}: arquivos menores com a mesma qualidade, suporte a resoluções e taxas de quadros mais altas.
  \item \textbf{Desvantagens}: codificação e decodificação exigem processador mais potente, alguns dispositivos antigos não são compatíveis.
  \item \textbf{Cenário de uso}: gravação de vídeo 4K, cenários com espaço de armazenamento limitado.
\end{itemize}

\subsubsection*{Apple ProRes:}

\begin{itemize}
  \item \textbf{Característica}: codificação profissional de vídeo desenvolvida pela Apple, compressão praticamente sem perdas.
  \item \textbf{Vantagens}: qualidade de imagem altíssima, grande margem para pós-produção, cores precisas.
  \item \textbf{Desvantagens}: arquivos enormes, exigem muito espaço de armazenamento, processamento requer equipamento profissional.
  \item \textbf{Cenário de uso}: produção de vídeo profissional, gravação de nível cinematográfico.
\end{itemize}

\subsubsection*{ProRes RAW:}

\begin{itemize}
  \item \textbf{Característica}: combina as vantagens dos formatos ProRes e RAW.
  \item \textbf{O que é RAW}: dados brutos do sensor, não processados, preservando o máximo de informação.
  \item \textbf{Vantagens}: margem de ajuste na pós-produção extremamente ampla, permite alterar exposição, balanço de branco e outros parâmetros sem perda de qualidade.
  \item \textbf{Desvantagens}: arquivos extremamente grandes, exigem software profissional de pós-produção.
  \item \textbf{Cenário de uso}: produção profissional de alto nível. A linha iPhone 15 Pro e posteriores suporta gravação de vídeo Apple ProRes (até 4K 60 fps); o iPhone 17 Pro vai além, suportando gravação direta de ProRes RAW no próprio aparelho.
\end{itemize}

\booksubsection{III. HDR: Tecnologia de Alto Alcance Dinâmico}

A tecnologia HDR (High Dynamic Range) preserva simultaneamente os detalhes das áreas mais claras e mais escuras da imagem.

\begin{itemize}
  \item \textbf{HDR10}: padrão HDR mais básico, boa compatibilidade.
  \item \textbf{HDR10+}: padrão promovido pela Samsung, suporta metadados dinâmicos.
  \item \textbf{Dolby Vision}: padrão HDR de topo desenvolvido pela Dolby, melhor qualidade de imagem, mas requer licenciamento.
  \item \textbf{HLG}: padrão HDR para broadcast, compatível com monitores comuns.
\end{itemize}

Ao gravar vídeo HDR, é necessário um monitor compatível com HDR para visualizar o efeito completo.

\booksubsection{IV. Espaço de Cor e Profundidade de Cor}

\subsubsection*{Espaço de Cor:}

\begin{itemize}
  \item \textbf{sRGB}: espaço de cor padrão, melhor compatibilidade.
  \item \textbf{Display P3}: padrão de gama ampla promovido pela Apple, cores mais ricas.
  \item \textbf{Rec.2020}: espaço de cor padrão para vídeo HDR, a gama de cores mais ampla.
\end{itemize}

\subsubsection*{Profundidade de Cor:}

\begin{itemize}
  \item \textbf{8-bit}: 16,7 milhões de cores, padrão de vídeo comum.
  \item \textbf{10-bit}: 1,07 bilhão de cores, padrão de vídeo profissional, transições de cor mais suaves.
  \item \textbf{12-bit}: 68 bilhões de cores, padrão profissional de alto nível.
\end{itemize}

\begin{tcbinfo}{Sugestões Práticas}
\begin{itemize}
  \item Gravação cotidiana: 1080p/4K + 30 fps + codificação H.265, equilibrando qualidade e tamanho de arquivo.
  \item Gravação de esportes: 1080p/4K + 60 fps, para garantir fluidez.
  \item Criação profissional: 4K + ProRes (se suportado), preservando máxima margem de pós-produção.
  \item Sugestão de armazenamento: vídeos de alto bitrate ocupam muito espaço; recomenda-se backup periódico no computador ou na nuvem.
\end{itemize}

\end{tcbinfo}

\begin{tcbpractice}{Explore: As Configurações de Vídeo do Seu Celular}
\begin{enumerate}
  \item Abra o app da câmera e mude para o modo de vídeo.
  \item Veja as opções de resolução e taxa de quadros e experimente diferentes combinações.
  \item Se houver suporte, verifique se há opções avançadas como HDR, ProRes, etc.
  \item Grave alguns trechos com configurações diferentes e compare a qualidade de imagem e o tamanho dos arquivos.
\end{enumerate}

\end{tcbpractice}

\booksubsection{V. Formato RAW vs. HEIF/JPEG}

\subsubsection*{O que é o formato RAW:}

O formato RAW são os “dados brutos” registrados pelo sensor da câmera — informação pura e original, sem qualquer compressão ou processamento. É como ingredientes frescos, que precisam ser “cozidos” na pós-produção para se tornarem o prato final. Como o celular tem armazenamento e poder de processamento limitados, por padrão ele gera fotos que já estão “prontas” (JPEG ou HEIF); o RAW precisa ser ativado especificamente.

\subsubsection*{Comparação RAW vs HEIF vs JPEG:}

\begin{itemize}
  \item \textbf{JPEG}: o formato mais antigo, melhor compatibilidade (qualquer dispositivo abre). Profundidade de cor de 8 bits, compressão com perdas, arquivo menor (2 a 5 MB por foto). Margem de pós-produção limitada; forçar o clareamento de sombras tende a gerar manchas de cor.
  \item \textbf{HEIF/HEIC}: novo formato promovido pela Apple. Com a mesma qualidade de imagem, o arquivo tem metade do tamanho do JPEG (1 a 3 MB por foto). Suporta profundidade de cor de 10 bits e canal alfa. A desvantagem é que alguns dispositivos antigos e o Windows precisam instalar plugin para abrir.
  \item \textbf{RAW}: dados brutos não comprimidos, arquivo maior (20 a 100 MB por foto). Profundidade de cor geralmente de 12 a 16 bits, máxima margem de pós-produção. A desvantagem é que exige software profissional (Lightroom, Snapseed, etc.) para abrir e processar; não pode ser compartilhado diretamente.
\end{itemize}

\subsubsection*{Por que fotógrafos profissionais fotografam em RAW:}

\begin{itemize}
  \item \textbf{Balanço de branco ajustável sem perdas na pós-produção}: no JPEG, se errar o balanço de branco (ex.: tom amarelado sob luz incandescente), corrigir depois degrada a qualidade; no RAW, é possível redefinir a temperatura de cor sem nenhuma perda.
  \item \textbf{Grande tolerância de exposição}: no JPEG, sombras subexpostas em 2 stops são basicamente pretas; ao clarear, só aparecem ruídos. O RAW consegue "resgatar" detalhes utilizáveis de sombras subexpostas em 3 a 4 stops. O mesmo vale para superexposição.
  \item \textbf{Detalhes de sombra e altas luzes mais ricos}: a faixa dinâmica do RAW é muito maior que a do JPEG; a margem para evitar que altas luzes estourem e sombras virem preto total é muito maior.
  \item \textbf{Nitidez, redução de ruído e correção de lente controláveis na pós-produção}: o processamento interno da câmera é uma "decisão definitiva"; no RAW, é possível ajustar repetidamente, sempre partindo dos dados originais.
\end{itemize}

\subsubsection*{Como ativar o RAW no celular:}

\begin{itemize}
  \item \textbf{iPhone ProRAW}: apenas modelos iPhone 12 Pro e superiores (linha Pro). Em "Ajustes Câmera Formatos", ative o "Apple ProRAW". Em seguida, aparecerá um botão "RAW" no topo da interface da câmera; toque para ativar. Arquivo de cerca de 25 MB por foto; recomenda-se ativar apenas quando necessário.
  \item \textbf{Android Motion Raw / Ultra RAW}: cada marca tem um nome diferente. Nos modelos flagship da Xiaomi/OPPO/vivo, no "Modo Profissional" ou "Mais" da câmera, há a opção RAW; na Huawei, ative o RAW no "Modo Profissional"; na linha Samsung Galaxy S, há a opção RAW no modo "Profissional". Tamanho de arquivo geralmente entre 20 e 50 MB.
  \item \textbf{Sugestão geral}: no dia a dia, não é necessário ativar o RAW — o arquivo é muito grande e ocupa espaço; ative apenas em situações importantes (paisagem, retrato, quando houver necessidade de pós-produção).
\end{itemize}

\booksubsection{VI. Fundamentos da Fotografia Computacional}

A fotografia móvel moderna já não pode prescindir da “fotografia computacional” — ou seja, gerar fotos com base em algoritmos, e não apenas na lente e no sensor. Entender esses princípios essenciais ajuda a compreender melhor o que o celular faz no instante em que o obturador é pressionado.

\subsubsection*{Síntese de múltiplos quadros:}

Quando você pressiona o obturador no celular, na verdade são capturadas várias fotos. O algoritmo alinha os vários quadros, pondera e os funde, obtendo uma imagem final de alta qualidade.

\begin{itemize}
  \item \textbf{Modo Noturno}: várias fotos subexpostas são capturadas em 1 a 3 segundos; alinhadas e sobrepostas — os ruídos se cancelam e o brilho se acumula. O resultado é uma imagem limpa mesmo em luz baixa.
  \item \textbf{HDR (High Dynamic Range)}: são capturadas separadamente fotos subexpostas, normais e superexpostas (três ou mais), e depois combinadas numa única foto com detalhes tanto nas altas luzes quanto nas sombras. Especialmente eficaz em cenas com contraluz.
  \item \textbf{Deep Fusion (Apple) / Photonic Engine (Motor Fotônico)}: voltado para cenas de luz média a baixa, analisa pixel a pixel vários quadros e os funde, com foco em melhorar textura e detalhes; detalhes como suéteres e cabelos ficam nitidamente mais nítidos.
\end{itemize}

\subsubsection*{Segmentação semântica:}

A IA não simplesmente aplica um processamento uniforme na imagem inteira; ela primeiro “entende” o que há na cena — céu, rosto, grama, construções, roupas — e então aplica estratégias de processamento distintas para cada região.

\begin{itemize}
  \item \textbf{Céu}: realce individual do azul, redução de brilho, para evitar superexposição.
  \item \textbf{Rosto}: identifica a região do tom de pele, clareia e suaviza a pele, sem mexer nos traços faciais e nos detalhes do cabelo.
  \item \textbf{Grama / Folhas}: realce da saturação do verde, aumento da nitidez dos detalhes.
  \item \textbf{Desfoque de fundo}: no modo retrato, a IA reconhece o contorno do tema principal, desfoca o fundo e mantém o tema nítido.
\end{itemize}

\subsubsection*{Aplicações de deep learning:}

\begin{itemize}
  \item \textbf{Reconhecimento de cena por IA}: a câmera identifica automaticamente cenas como "comida", "animal de estimação", "noturno", "documento" e ajusta os parâmetros de forma otimizada. Por exemplo, o modo comida realça tons quentes e saturação, deixando o alimento mais apetitoso.
  \item \textbf{Desfoque de retrato por IA}: a lente do celular fisicamente não consegue a profundidade de campo rasa de uma DSLR; ela depende de modelos de deep learning para prever as bordas do tema e então desfocar o fundo. Os algoritmos mais recentes já processam corretamente detalhes como fios de cabelo e bordas de vidro.
  \item \textbf{Remoção por IA / Expansão por IA}: selecione pessoas ou objetos indesejados na imagem; a IA "completa" as áreas ocluídas com base no contexto. É uma aplicação da IA generativa no campo da imagem.
\end{itemize}

\subsubsection*{Fotografia computacional vs. fotografia óptica:}

\begin{itemize}
  \item \textbf{Vantagens da fotografia computacional}: tamanho reduzido também consegue boa qualidade de imagem; cenas extremas como baixa luz e contraluz também produzem boas fotos; o usuário comum não precisa conhecer fotografia para tirar fotos "bonitas".
  \item \textbf{Desvantagens da fotografia computacional}: a imagem às vezes fica "superprocessada" — nitidez excessiva, distorção de cor, cor do céu não natural; a síntese de múltiplos quadros deixa rastros quando o tema se move; os detalhes são "calculados", não registrados de verdade.
  \item \textbf{Vantagens da fotografia óptica}: sensor grande + lente de grande abertura capturam a luz de verdade, qualidade de imagem natural, grande margem de pós-produção; não produz artefatos porque o algoritmo "não inventa".
  \item \textbf{Tendência}: os celulares flagship modernos combinam "óptica + computação" — usam sensor grande para obter informação real, e algoritmos para otimizar a apresentação. As duas abordagens não são opostas.
\end{itemize}

\booksubsection{VII. Tipos de Lentes e Distâncias Focais}

Os celulares flagship modernos geralmente têm três ou quatro lentes, cada uma correspondendo a uma distância focal diferente. Entender a distância focal permite escolher a lente certa em cada situação.

\subsubsection*{Lentes comuns e distâncias focais equivalentes:}

\begin{itemize}
  \item \textbf{Ultra grande-angular (equivalente a 13-24 mm)}: campo de visão mais amplo, ideal para paisagens, arquitetura e espaços internos. Capta a grande cena que você vê diante dos olhos, mas as bordas apresentam distorção (linhas retas se curvam); em retratos, o rosto é esticado. Bom para "grandes cenas", não para "close-up".
  \item \textbf{Câmera principal (equivalente a 24-28 mm)}: a lente mais usada no dia a dia, com a melhor qualidade de imagem e maior abertura. Equilibra campo de visão e qualidade; funciona para pessoas e paisagens. O "1x" da maioria dos celulares é a câmera principal.
  \item \textbf{Teleobjetiva (equivalente a 70-200 mm)}: para fotografar paisagens distantes e retratos em close-up. A teleobjetiva comprime a sensação de espaço (montanhas distantes e pessoas próximas parecem mais perto); em retratos, produz proporções faciais mais naturais e desfoque de fundo.
  \item \textbf{Teleobjetiva periscópica (equivalente a 200 mm+)}: por meio de um design de caminho óptico reflexivo, alcança distâncias focais maiores dentro da espessura do celular. Permite capturar detalhes nítidos a grande distância; ideal para shows, eventos esportivos e vida selvagem. Presente apenas nos modelos flagship.
\end{itemize}

\subsubsection*{Compromisso entre tamanho do sensor e tamanho do pixel:}

O sensor é a “base” que determina quanta luz pode ser recebida. No jargão do setor, “uma base maior esmaga a menor” — com o mesmo número de pixels, quanto maior o sensor, maior cada pixel individual, mais luz entra e melhor a qualidade de imagem.

\begin{itemize}
  \item \textbf{Sensor grande}: grande entrada de luz, boa qualidade de imagem noturna, ampla faixa dinâmica. A desvantagem é que o celular precisa ser mais grosso (módulo da lente saliente), e o custo é mais alto.
  \item \textbf{Alta resolução com sensor pequeno}: com boa luz do dia, a alta resolução produz fotos mais nítidas, facilitando o recorte posterior. Mas em luz baixa, cada pixel recebe luz insuficiente, gerando muito ruído.
  \item \textbf{Escolha dos flagships atuais}: a câmera principal costuma ser de 1/1,3 polegada a 1 polegada (alguns modelos já usam sensor de 1 polegada), com cerca de 50 megapixels — é o ponto de equilíbrio entre qualidade de imagem e volume.
\end{itemize}

\subsubsection*{Tecnologia de fusão de pixels (Pixel Binning):}

Para conciliar alta resolução de dia com pixels grandes à noite, os fabricantes usam a “fusão de pixels” — vários pixels pequenos são combinados num único pixel grande de saída.

\begin{itemize}
  \item \textbf{Fusão de quatro pixels (Quad Bayer)}: sensor de 50 megapixels, gerando foto de 12,5 megapixels. Cada pixel de saída é formado pela fusão de 4 pixels adjacentes; a entrada de luz aumenta 4 vezes, com qualidade de imagem mais limpa.
  \item \textbf{Fusão de nove pixels (Nonacell)}: usado em sensores acima de 100 megapixels, gerando cerca de 11 a 12 megapixels. Mesmo princípio, com entrada de luz ainda maior.
  \item \textbf{Efeito prático}: no modo de foto padrão, você recebe a foto com pixels fundidos — qualidade boa, arquivo pequeno. Quando precisar de detalhes extremos (como fotografar documentos ou paisagens distantes), é possível mudar manualmente para o "Modo de Alta Resolução" para obter a resolução total.
\end{itemize}

\booksubsection{VIII. Modos de Captura Especiais}

Além da fotografia cotidiana, o celular tem alguns modos especiais que, bem usados, capturam imagens que normalmente não seriam possíveis.

\subsubsection*{Modo Panorama:}

\begin{itemize}
  \item Ao fotografar, mova o celular lenta e suavemente na horizontal conforme a seta indicativa; não vá rápido — se for rápido, aparecem falhas na costura.
  \item Mantenha o celular estável na vertical, movendo apenas na horizontal. Se balançar para cima e para baixo, a linha do horizonte vira uma onda.
  \item Cuidado com pessoas ou carros em movimento — na costura, podem aparecer "meio corpo" ou "fantasmas".
  \item O modo panorama também pode ser usado na vertical — fotografar na vertical com costura horizontal é ideal para edifícios altos ou cachoeiras.
\end{itemize}

\subsubsection*{Câmera Lenta e Time-Lapse:}

\begin{itemize}
  \item \textbf{Câmera lenta}: grava com alta taxa de quadros (120 fps ou 240 fps); na reprodução, vira câmera lenta. Ideal para gotas d'água, momentos esportivos, saltos de animais. A luz precisa ser abundante; no modo câmera lenta, o tempo de entrada de luz é curto, e a qualidade em baixa luz é ruim.
  \item \textbf{Time-lapse}: captura fotos em intervalos e as combina num vídeo acelerado. Ideal para fluxo de carros, movimento de nuvens, pôr do sol, abertura de flores. O tempo de captura varia de alguns minutos a várias horas; é imprescindível usar tripé para fixar o celular; na mão, com certeza fica borrado.
  \item \textbf{Dica}: para time-lapse de nuvens, monte o celular, grave por 15 a 30 minutos e obtenha um vídeo de time-lapse de cerca de 10 segundos — o efeito é impactante.
\end{itemize}

\subsubsection*{Modo Estrelas / Astrofotografia:}

\begin{itemize}
  \item É preciso se afastar da poluição luminosa urbana; procure áreas rurais ou montanhosas totalmente sem interferência de postes de luz. Na cidade, não se fotografa a Via Láctea.
  \item É imprescindível usar tripé; na mão, com certeza fica borrado. O celular precisa ficar firmemente apoiado por vários segundos a dezenas de segundos.
  \item O tempo precisa estar limpo e sem nuvens; a lua não deve estar muito grande (lua cheia ofusca a luz das estrelas). As noites escuras do início ou fim do mês lunar são as mais adequadas.
  \item O "Modo Noite" do iPhone reconhece automaticamente o céu estrelado e prolonga o tempo de exposição (até 30 segundos); o "Modo Estrelas" ou "Modo Astrofotografia" dos flagships Android tem efeito similar. É preciso paciência; depois de pressionar o obturador, não mexa no celular.
\end{itemize}

\subsubsection*{Modo Macro:}

\begin{itemize}
  \item Usado para fotografar flores, insetos, joias, close-ups de texto e outros objetos pequenos. Permite focar a distâncias de 1 a 2 cm.
  \item Alguns celulares têm lente macro dedicada (geralmente sensor pequeno de 2 ou 5 megapixels), com qualidade mediana; os flagships usam a lente ultra grande-angular para também fazer macro, com qualidade melhor.
  \item Ao fotografar de muito perto, a lente fica muito próxima do tema e pode bloquear a luz; se necessário, use um abajur ou lanterna para iluminar.
  \item A distância de foco é muito curta e a profundidade de campo é muito rasa; mover o tema alguns milímetros para frente ou para trás já pode desfocar. A mão precisa estar firme; de preferência, mantenha o tema imóvel.
\end{itemize}

\begin{tcbinfo}{Sugestões Práticas de Fotografia Estática}
\begin{itemize}
  \item Fotografia cotidiana: HEIF/JPEG padrão é suficiente; arquivo pequeno, qualidade bastante boa.
  \item Cenas importantes (paisagem, retrato, necessidade de pós-produção): ative o RAW, preservando máxima margem de pós-produção.
  \item Escolha a lente certa: paisagem com ultra grande-angular, retrato com teleobjetiva, dia a dia com a câmera principal.
  \item Noturna e estrelas: imprescindível fixar o celular; ative o modo noturno ou modo astrofotografia.
  \item Celulares com sensor grande (como modelos de 1 polegada) têm teto de qualidade de imagem mais alto, mas os hábitos de operação continuam sendo o fator decisivo.
\end{itemize}

\end{tcbinfo}

\begin{tcbpractice}{Mãos à Obra: Conheça a Capacidade da Sua Câmera}
\begin{enumerate}
  \item Verifique se seu celular suporta captura RAW (Ajustes Câmera Formatos); tente ativar e fotografar, comparando o tamanho do arquivo JPEG e RAW da mesma cena.
  \item Mude para as três distâncias focais — ultra grande-angular, câmera principal e teleobjetiva — e fotografe a mesma paisagem distante com cada uma, percebendo a diferença de perspectiva e de qualidade de imagem.
  \item Encontre uma cena adequada para modo noturno ou modo estrelas; fixe o celular e fotografe uma foto noturna.
  \item Experimente o modo panorama na horizontal para uma grande cena, cuidando para mover o celular de forma estável.
\end{enumerate}

\end{tcbpractice}
